% Transcription par lecture visuelle directe, sans OCR.
% Statut : transcrit ; relecture indépendante à effectuer.
% Source : PDF original, vue 109, côté gauche, page manuscrite 218.
\chapter{Préfecture de M. Franchet}
\begin{center}
Suspension des assemblées générales.
\end{center}

\section*{Son installation}
% Mention marginale de la page manuscrite 218.
La position de la Congrégation devenait de plus en plus critique ; heureusement son nouveau préfet avait toutes les qualités nécessaires pour qu'il pût se placer à la hauteur des devoirs qu'une charge rendue bien pesante par les circonstances allait lui imposer. Jusqu'à ce jour chaque préfet de la Congrégation avait reçu de la Providence sa mission spéciale. Nous avons vu la Congrégation se fonder sous l'autorité de M. Ripoud, sous celle de M. Chaulet s'accroître et se fortifier. M. Philpin a été destiné à la soutenir au milieu des épreuves. Avec M. Franchet elle va entrer dans la voie des persécutions, son préfet y marchera à sa tête à la suite du souverain pontife, et s'il ne lui est pas donné d'obtenir la couronne du martyre, il aura du moins acquis un droit incontestable à celle des confesseurs. Dieu l'avait façonné tout exprès pour le ministère qui lui était réservé ; M. Franchet était un homme de foi et savait au besoin être un homme d'action. Il avait le jugement droit, le coup d'œil prompt et juste et une volonté forte. On pouvait compter également sur sa prudence pour prévoir le danger et sur sa fermeté pour le mépriser et lui opposer lorsqu'il était inévitable une constance à toute épreuve. Ceux qui le connaissaient intimement pouvaient déjà pressentir que dans les conseils du souverain maître, il avait été choisi pour servir efficacement la cause de sa religion. Dieu
% Source : PDF original, vue 109, côté droit, page manuscrite 219.
avait ses desseins sur lui, il était destiné aux honneurs, mais c'est par la voie des persécutions qu'il devait y parvenir.

La Congrégation était encore debout au moment de son installation, rien n'était changé à sa manière d'être ; les assemblées avaient lieu comme à l'ordinaire, les sections exerçaient toujours les mêmes bonnes œuvres ; il faut toutefois en excepter la visite des prisonniers que déjà la section des prisons et hôpitaux avait été forcée d'abandonner, mais elle s'en dédommageait en donnant plus d'activité aux soins qu'elle prenait des malades de l'Hôtel-Dieu.

Ces restes de paix et de tranquillité dont nous jouissions encore ne pouvaient pas se prolonger bien longtemps ; nous ne pouvions guère espérer qu'on nous refusât toute espèce de participation à la persécution élevée contre le chef de l'Église. La guerre que Napoléon avait déclarée à son autorité se poursuivait avec fureur. L'insensé ne comprenait pas que c'était en vain qu'il méditait de détruire une puissance établie par Dieu lui-même et qui a pour appui la force de son bras. Le Seigneur a parlé et sa parole ne passera point, mais les passions des hommes dans leur délire ne veulent pas entendre ce langage.

L'empereur, dans l'espoir d'enlever au pape le droit dont il est en possession d'instituer les évêques, convoqua à Paris sous sa main despotique un concile composé de tous les évêques de France et d'Italie. Les actes trop nombreux de foiblesse que l'Église avait eu à déplorer, lui faisaient préjuger qu'il obtiendrait au moyen de ce concile un succès complet.
% Source : PDF original, vue 110, côté gauche, page manuscrite 220.
\section*{Concile national}
% Mention marginale de la page manuscrite 220.
Les évêques s'assemblent dans la basilique de Paris. Mgr de Boullogne, évêque de Troyes, prononce en présence de tous les pères, un discours éloquent sur l'autorité de l'Église, dans lequel il célèbre en même temps celle de son chef. Ainsi s'ouvrit la première session. Le concile rendit deux décrets solennels ; par le premier, il reconnaît la primauté de juridiction et d'honneur que le Saint-Père exerce dans l'Église universelle. Par le second, il déclare se conformer entièrement à la doctrine du concile de Trente, doctrine que l'on suivait en France comme en Italie, mais à laquelle l'Église de France n'avait jamais donné un assentiment aussi formel. Là se bornèrent les travaux de ce concile convoqué et ouvert avec tant d'éclat. À la vérité, on tint encore plusieurs congrégations ; (vous savez qu'on appelle ainsi les conférences préliminaires dans lesquelles les pères d'un concile préparent d'avance les canons ou décrets qu'ils se proposent de promulguer dans les sessions.) Mais la majorité s'étant montrée entièrement opposée aux volontés schismatiques de l'empereur, celui-ci dans sa colère fit arrêter plusieurs évêques et se hâta de congédier les autres.

Ainsi encore une fois l'iniquité s'est menti à elle-même et s'est embarrassée dans ses propres filets. Ceux qui comme autrefois Balaam avaient été appelés pour maudire dirigés, comme lui par l'Esprit-Saint, n'ont pu que bénir. Reconnaissons donc ici la présence toujours protectrice de celui qui a dit : Et voilà que je suis avec vous jusqu'à la consommation des siècles.

Le résultat du concile avait réjoui les âmes fidèles, mais il
% Source : PDF original, vue 110, côté droit, page manuscrite 221.
avait comme cela devait être augmenté l'exaspération de Napoléon. Il s'irritait de plus en plus à mesure qu'il voyait de nouveaux obstacles s'élever contre sa volonté. Dans sa fureur il cherchait constamment les moyens d'ôter au pape tous les appuis qui paraissaient encore soutenir sa puissance. Nous l'avons déjà vu renverser l'ordre des Pères de la Foi ; maintenant c'est aux Congrégations qu'il va s'en prendre, car il ne doute pas que ces sociétés ne soient composées de chrétiens tout dévoués au souverain pontife. Voilà bien les passions humaines, elles produisent toujours les mêmes effets. Autrefois Hérode tremblait sur son trône à la nouvelle de la naissance d'un enfant, et maintenant un illustre conquérant tonne de tous ses foudres contre quelques jeunes gens désarmés qui dans plusieurs villes de France se réunissent pour prier Dieu et pour soulager leur prochain.

Les Congrégations de Paris et de Bordeaux avaient beaucoup plus de notoriété que les autres ; les premiers coups durent être dirigés contre elles. Ces deux sociétés sont dissoutes par ordre de l'empereur. Vers la même époque celle de Paris avait fait une perte irréparable ; Dieu avait appelé à lui le vénérable P. Delpuy, son directeur ; seul il était l'âme de cette Congrégation, il en tenait dans sa main tous les fils, sa mort achevait nécessairement de la dissoudre ; s'il eût vécu, il aurait pu sans doute concentrer son petit noyau autour de lui et lui conserver à l'ombre du mystère une existence secrète ; mais toutes ces espérances se trouvèrent anéanties car personne autre n'aurait pu les réaliser.

Après ce coup d'essai, pour atteindre plus sûrement toutes
% Source : PDF original, vue 111, côté gauche, page manuscrite 222.
les sociétés de ce genre qui de tous côtés s'étaient élevées dans la France, ou du moins pour ne leur laisser d'autre alternative que de s'enchaîner à son char ou de se dissoudre, l'empereur fit insérer dans le code pénal qui se discutait alors, un article par lequel il était défendu sous diverses peines, de faire partie d'une association ayant pour but la politique, la littérature ou la religion, si elle était composée de plus de vingt membres, à moins que cette association ne fût autorisée par l'empereur. Des ordres sévères furent envoyés dans les départements pour mettre de suite cette mesure à exécution. Ces ordres parvinrent à M. Maillocheau, commissaire général de police à Lyon ; on lui mandait en même temps qu'il existait une Congrégation consacrée à la Ste Vierge et qu'il fallait la soumettre à la nouvelle législation.

\section*{Décret contre les Congrégations}
% Mention marginale de la page manuscrite 222 ; placée au début de l'application locale.
Voilà M. Maillocheau à la piste de la Congrégation. Il s'adresse à tous les curés de la ville pour savoir où il pourra la découvrir. M. le curé de St Nizier répond qu'il ne connaît dans sa paroisse d'autre confrérie de la Ste Vierge que celle de N. Dame de grâces, presque entièrement composée de femmes. Tous les autres curés accusent l'existence de la confrérie du Rosaire qui était tout à fait dans le même cas. Le pauvre commissaire de police était désappointé et la Congrégation, à l'ombre du mystère toutefois assez transparent dont elle était enveloppée, n'en continuait pas moins à agir comme par le passé, en attendant qu'on fût parvenu à la trouver, tout en n'ignorant pas néanmoins les démarches qu'on faisait pour la découvrir.

M. Maillocheau avait d'abord frappé à faux ; mais il
% Source : PDF original, vue 111, côté droit, page manuscrite 223.
finit enfin par rencontrer juste, il s'adressa à l'archevêché. Là il trouva M. Bochard qui s'empressa de lui donner connaissance de l'existence de notre société ; le commissaire de police lui demanda alors communication des règlements et la liste des noms des personnes qui en faisaient partie. L'assurant que s'il ne s'y trouvait rien de contraire aux intentions de l'empereur, il se faisait fort d'obtenir l'autorisation du gouvernement. M. Bochard qui ne considérait notre société que comme une simple société de bonnes œuvres ne vit dans toutes ces exigences que l'avantage de nous obtenir plus de liberté et une sécurité complète, il ne mit pas en doute que la Congrégation envisagerait cette affaire du même œil que lui, et s'empresserait à se soumettre à ce qu'on demandait d'elle. Dans cette persuasion, il crut devoir donner au commissaire de police l'assurance qu'on le satisferait sur tous les points.

Il fit ensuite part à M. Franchet notre préfet des demandes qu'il s'était chargé de nous transmettre ; il put s'apercevoir dès lors qu'il pourrait bien se faire qu'il se fût trop avancé. M. Franchet lui répondit qu'avant de donner son consentement à une mesure aussi grave, il lui fallait nécessairement le temps de la réflexion, qu'il allait en conférer avec les assistants et qu'ensuite il lui rendrait réponse.

Effectivement il convoqua aussitôt les assistants, leur adjoignit quelques-uns des membres les plus au courant des affaires de la Congrégation et leur fit part que l'orage auquel nous nous attendions depuis longtemps, était prêt à éclater sur nous. La question fut mûrement examinée sur toutes ses faces et les avis furent unanimes ; on ne voulait consentir à aucun prix à se mettre sous la main d'un gouvernement qui avait alors
% Source : PDF original, vue 112, côté gauche, page manuscrite 224.
des tendances évidemment schismatiques, et qui ne nous flattait vraisemblablement que pour nous entraîner de son côté dans sa lutte avec le souverain pontife, sauf à nous perdre plus sûrement plus tard s'il ne pouvait pas y réussir. On prit donc la résolution inébranlable de repousser de quelque part qu'elle vînt, toute tentative de mêler le gouvernement directement ou indirectement à nos affaires ; on décida à plus forte raison de refuser toutes communications de listes ; cependant pour avoir l'air d'accorder quelque chose, comme on n'y vit pas un grand inconvénient on se décida à faire connaître le règlement.

Il ne restait plus qu'à décider par qui on ferait passer notre réponse ; c'était à M. Bochard qu'on s'était adressé ; il eût semblé naturel de le prier de s'en charger ; mais il nous avait déjà compromis, il pouvait nous compromettre plus gravement encore. M. Franchet s'offrit pour aller lui-même attacher le grelot ; cette offre fut acceptée avec reconnaissance ; nos intérêts ne pouvaient être placés en meilleures mains, car il possédait toute la présence d'esprit et la fermeté d'âme qu'exigeait une semblable mission.

Nous appelâmes à notre secours dans cette circonstance suivant notre usage notre auguste patronne. Comme toujours elle eut soin de nous, et grâce à sa puissante protection la Congrégation put résister à cette rude épreuve.

\section*{Conférence avec M. Maillocheau}
% Mention marginale des pages manuscrites 224 et 226 ; placée au début de l'entretien.
M. Franchet se munit d'une copie du règlement, dit son Veni Creator, et avec ce sang-froid qui ne l'abandonnait jamais, se présenta chez M. le commissaire général de police. Monseigneur, lui dit-il, vous avez témoigné à M. Bochard, vicaire général, le désir de connaître une petite société qui existe à Lyon,
% Source : PDF original, vue 112, côté droit, page manuscrite 225.
si vous voulez le permettre, je viens pour vous donner les renseignements que vous désirez, je puis vous en parler avec connaissance de cause car j'en fais partie. Cette société se compose de jeunes gens, elle n'a aucun but politique, elle borne tous ses soins à venir au secours des malheureux ; si vous jugez à propos de prendre connaissance de ses règles, je pourrai vous les communiquer.

M. Maillocheau reçut M. Franchet avec la plus grande politesse, il commença par donner de grands éloges au dévouement que ces jeunes gens montraient pour la cause de l'humanité ; mais ajouta-t-il, une société de ce genre ne peut agir dans l'ombre ; il faut qu'elle se fasse connaître : ainsi outre ses règlements, vous me remettrez la liste de ceux qui la composent et vous pouvez compter sur la protection du gouvernement. — Ce que vous me demandez là, Monseigneur, me paraît très difficile à concilier avec la nature de notre société ; elle tient essentiellement à opérer dans le secret le peu de bien dont elle est capable ; il serait donc tout à fait contraire à sa manière d'être de divulguer les noms de ceux qui en font partie ; vous voyez que pour mon compte je n'ai pas craint de vous dire que j'en suis, mais il y a parmi nous un grand nombre de jeunes gens que la publicité contrarierait beaucoup ; les uns sont en la puissance de leurs parents, les autres dépendent de chefs de commerce chez lesquels ils sont commis ; la plupart d'entre eux trouveraient fort désagréable que leurs rapports de bonnes œuvres fussent ainsi dévoilés. Ils n'ont consenti à faire partie de cette société précisément qu'à cause de l'incognito qu'elle garde ; vous comprenez qu'ils auraient lieu de se plaindre si aujourd'hui cela se passait différemment. — Eh bien, Monsieur, s'il en est ainsi demandez
% Source : PDF original, vue 113, côté gauche, page manuscrite 226.
donc leur agrément pour me faire cette communication ; elle est indispensable ; le gouvernement l'exige, et si vous vous y refusez, je ne pourrais pas obtenir que votre société fût tolérée ; il serait cependant bien fâcheux qu'une société comme la vôtre aussi philanthropique, fût obligée de se dissoudre. Prenez donc l'avis de vos Messieurs, je suis convaincu qu'ils ne se refuseront pas à ce que je leur demande. — Je prends, Monseigneur, l'engagement de soumettre aux membres de notre société toutes les observations que vous me faites ; mais je ne vous dissimule pas que j'ai l'intime conviction qu'ils se résigneront à se séparer, plutôt que de se soumettre à un système de publicité qui contrarierait toutes leurs idées. — Malgré tout ce que vous me dites, j'espère encore qu'il n'en sera pas ainsi ; voyez-les toujours. — Ainsi se termina cette conférence. M. Franchet se retira sans laisser la copie du règlement ; il se hâta d'aller faire part du résultat de cette démarche officielle à ceux de nos amis avec lesquels elle avait été convenue. Il jugea ensuite convenable de convoquer sa section, d'y joindre les chefs de section et quelques autres membres influents ; c'était comme une assemblée préparatoire pour délibérer sur les grandes mesures à prendre dans une circonstance aussi critique. Elle formait en nombre un peu plus du tiers de celui des membres de la Congrégation. Le préfet exposa à cette réunion la gravité des dangers dont la Congrégation était menacée et ne dissimula pas que quelque parti qu'on adoptât, la position demeurerait toujours hérissée de difficultés sans nombre. Il demanda à connaître l'opinion individuelle de chacun de ceux qui étaient présents sur la conduite que nous devions tenir.

% Source : PDF original, vue 113, côté droit, page manuscrite 227.
Tous à l'unanimité s'arrêtèrent à cette pensée que le plus grand péril auquel nous pouvions être appelés, serait très certainement celui qui résulterait de notre consentement à nous placer sous la griffe de la police en lui livrant nos noms ; qu'ainsi il valait mieux tout risquer plutôt que de donner les mains à une semblable mesure. En conséquence il fut décidé qu'il fallait annoncer ouvertement la dissolution de la Congrégation, suspendre effectivement les assemblées générales, mais tout en rentrant ainsi dans l'incognito le plus absolu pour laisser passer l'orage, n'en continuer pas moins de se réunir en sections dans le plus profond secret et faire isolément quant à l'action seulement, mais toujours dans un esprit d'union, toutes les bonnes œuvres qui par leur nature ne pouvaient pas déceler l'existence d'une société. Le préfet après avoir ainsi pressenti la manière dont la Congrégation envisagerait cet événement, alla trouver M. Bochard et lui fit part de la conférence qui avait eu lieu entre le commissaire général de police et lui et ne lui dissimula pas que son opinion personnelle était que nous ne pouvions pas sans commettre une imprudence vraiment impardonnable consentir à ce qu'on demandait de nous ; du reste, ajouta-t-il, il ne reste plus maintenant qu'à consulter la Congrégation. Il faudra lui mettre sous les yeux soit les avantages soit les inconvénients de ce qu'on propose. Lorsque chaque membre aura pu ainsi apprécier le pour et le contre, il se décidera avec connaissance de cause et le vote de la majorité fixera le sort de la Congrégation.

L'assemblée fut donc convoquée ; on eut soin d'exciter l'intérêt en annonçant qu'elle avait pour but des communications importantes. Déjà grand nombre de membres savaient ce dont
% Source : PDF original, vue 114, côté gauche, page manuscrite 228.
il était question, ils s'y rendirent donc avec empressement, il n'y avait pas encore eu d'assemblée aussi nombreuse. 46 congréganistes répondirent à l'appel. Après la lecture du procès-verbal M. Bochard prit le premier la parole ; il exposa dans un long discours et avec tout l'art d'une éloquence mielleuse tous les avantages que la Congrégation pourrait espérer en satisfaisant à ce que le gouvernement demandait d'elle. Elle allait trouver dans la bienveillance des autorités une protection puissante, dont le résultat serait pour elle une existence assurée. Il s'efforça de démontrer que la mesure impérieusement exigée de donner la liste de nos noms, ne pouvait pas présenter le moindre danger.

M. le directeur fut entendu avec respect, au milieu du plus profond silence et de l'attention la mieux soutenue. Nul doute qu'en terminant sa longue harangue, il n'eût la conviction d'avoir porté la persuasion dans tous les esprits.

M. le préfet prit ensuite la parole. Sans contester qu'il ne fût absolument possible que la Congrégation pût rencontrer quelques avantages dans une démarche qui avait pour but d'assurer sa sécurité, avantages qui ne lui paraissaient cependant pas démontrés avec une certaine évidence, M. Franchet insista avec beaucoup de force sur la position fausse dans laquelle cette concession de notre part allait nous placer ; il ne lui fut pas difficile de faire apercevoir aux moins clairvoyants dans quelle voie fatale et périlleuse nous nous laisserions entraîner si nous consentions à donner à la Congrégation une publicité si contraire à ses règles et à son esprit et qui dans les circonstances délicates où se trouvait l'Église pouvait l'exposer
% Source : PDF original, vue 114, côté droit, page manuscrite 229.
peut-être plus tard à se ranger parmi les ennemis du St Siège, dans la crainte de déplaire à ceux sous l'égide desquels elle se serait placée, ou si notre foi nous préservait de ce malheur, à nous trouver plus directement en butte à la persécution.

Lorsque M. Franchet eut cessé de parler, personne n'ouvrit la bouche pour demander la parole. L'assemblée tout entière attendait dans un silence qui avait quelque chose d'imposant et de solennel, le moment où elle allait faire connaître sa détermination. Le préfet posa ainsi la question : Que ceux qui sont d'avis de remettre au commissaire général de police la liste des membres de la Congrégation, se lèvent ? Un seul membre est debout. 45 restent assis. Seconde question : Que ceux qui sont d'avis de dissoudre la Congrégation plutôt que de jamais consentir à livrer les noms des congréganistes se lèvent ? 45 membres sont debout, un seul reste assis. Le résultat de cette délibération qui offrait pourtant une unanimité bien plus compacte qu'on ne l'aurait imaginé, était attendu par tout le monde si on en excepte toutefois M. Bochard. Aussi en fut-il stupéfait, on le vit pâlir, se troubler, en un mot paraître entièrement déconcerté. L'assemblée conservait son calme et son attitude également ferme et silencieuse ; ce fut ainsi qu'elle entendit son préfet proclamer que par suite de la décision qu'elle venait de prendre, la Congrégation pouvait dès ce jour se considérer comme dissoute ; cela n'empêcha pas de continuer à tenir l'assemblée jusqu'à la fin suivant ses us et coutumes sans en omettre le moindre détail. Lorsqu'elle eut été terminée par la prière de clôture, M. Bochard qui paraissait piqué
% Source : PDF original, vue 115, côté gauche, page manuscrite 230.
du peu d'effet qu'avaient produit ses exhortations en faveur des exigences de la police, se hâta de se retirer. Aussitôt après son départ, on fit circuler parmi tous les membres l'avis qu'il ne fallait considérer la dissolution de la société qui venait d'être votée que comme la suspension des assemblées générales et de tous les rapports extérieurs, mais que nous n'en resterions pas moins comme par le passé unis par le désir de procurer la gloire de Dieu et de nous entraider dans la grande affaire de notre sanctification réciproque, que nous continuerions à nous occuper de bonnes œuvres en redoublant toutefois de prudence, qu'ainsi il ne faudrait pas manquer de se rendre exactement comme à l'ordinaire aux assemblées de section ; tout cela fut parfaitement compris et fidèlement exécuté.

\section*{Dissolution apparente de la Congrégation}
% Mention marginale de la page manuscrite 230 ; placée au début des suites de la décision.
Il restait au préfet un devoir à remplir pour la solution de cette grave affaire, c'était d'aller faire connaître à M. le commissaire général de police le refus formel de la Congrégation d'obtempérer à sa demande. Celui-ci ne s'attendait pas, dans un temps où tout pliait devant la volonté du despote et celle des agents de sa tyrannie, à rencontrer tant de fermeté dans une société qui en elle-même paraissait bien peu importante. Il en fut tout étonné. C'est donc à dire, Messieurs, répondit-il à M. Franchet, que vous ressemblez à des lapins ; dès qu'on veut vous voir de près, vous rentrez dans votre trou… M. Franchet lui répliqua avec calme que ce résultat ne l'avait nullement surpris, qu'il l'avait bien prévu et le lui avait annoncé d'avance, qu'il était bien facile de comprendre que des jeunes gens qui se plaisaient à employer leurs loisirs à des bonnes œuvres, ne se souciassent cependant pas, pour faire le bien,
% Source : PDF original, vue 115, côté droit, page manuscrite 231.
de s'exposer à être, pour ainsi dire, affichés et montrés au doigt par tout le monde, qu'ils n'avaient pas hésité un instant, qu'ils aimaient mille fois mieux y renoncer pour toujours. — Ce n'est pas ainsi que je l'entends, répondit vivement M. Maillocheau ; si vous cessez de faire ces œuvres philanthropiques en société, j'espère au moins que chacun de vous ne cessera pas de s'en occuper individuellement. — Je crois pouvoir vous en donner l'assurance au nom de tous, du moins autant que cela nous sera possible, dit à son tour M. Franchet, bien satisfait de rencontrer dans cette dernière parole un excellent moyen de justifier au besoin les œuvres que la Congrégation se proposait de continuer à opérer par les sections. Si nous enfouissions ainsi notre union dans l'obscurité du mystère, ce n'était que pour la rendre plus inaltérable en la mettant à l'abri de tout danger ; le projet de dissoudre réellement la Congrégation n'était venu à la pensée de personne ; elle passa cependant pour l'être tout de bon aux yeux de M. Bochard ; dès ce moment tous rapports avec lui cessèrent entièrement.

Je me hâte d'ajouter pour votre édification, que dans cette position si délicate où les événements venaient de nous placer, dans cette espèce de guerre ouverte avec le délégué de l'autorité spirituelle, la Congrégation forcée de déposer le guide qui aurait dû être son père et son conseil ne demeura pas entièrement abandonnée à son propre mouvement. Toujours soumise aux supérieurs même au moment où elle aurait pu avoir quelque sujet de s'en plaindre, elle fut assez heureuse pour que ses chefs rencontrassent dans la personne de M. Caburra, saint prêtre qui avait appartenu à la société des Pères de la Foi, et qui avait alors des pouvoirs de vicaire général, un homme prudent et éclairé bien capable
% Source : PDF original, vue 116, côté gauche, page manuscrite 232.
de la diriger par ses conseils. La Providence semblait l'avoir placé là tout exprès pour servir de boussole à la Congrégation, l'aider à traverser en assurance le passage étroit dont elle ne devait pas s'écarter, sous peine d'aller échouer contre un des deux écueils qu'elle avait à redouter en même temps, d'un côté celui de manquer à la soumission et à la déférence qu'elle devait aux pasteurs du diocèse, de l'autre celui de s'exposer à être placé hors de l'obéissance et de la fidélité à toute épreuve qu'elle était bien résolue de conserver à jamais envers le chef de l'Église. M. Caburra fut donc alors d'une bien grande utilité à la Congrégation.

\section*{Choix du directeur des Congrégations}
% Mention marginale de la page manuscrite 232 ; placée au début du rôle officiel du directeur.
Par sa position officielle, M. Caburra, premier vicaire général, était obligé à de grands ménagements. Cela l'empêchait de prendre ouvertement le parti de la Congrégation, mais il l'aimait néanmoins beaucoup et la favorisait par dessous main de tout son pouvoir ; cette protection ne lui a jamais manqué jusqu'au moment où des temps plus heureux ont enfin permis à notre société de se replacer complètement sous l'autorité spirituelle.

Cette triste époque était partout fatale aux Congrégations. Nous avons déjà vu succomber celle de Bordeaux et celle de Paris. Peu de temps après, celle de Grenoble et le noyau formé à St Étienne eurent le même sort. Les amis de Chambéry disparurent aussi dans la tempête.

La Congrégation de la Ste Famille établie à Lyon pour les hommes mariés, sans être dissoute régulièrement, l'était de fait. Non seulement elle n'avait plus d'assemblées générales, mais ses sections elles-mêmes avaient cessé de s'assembler ; on n'y indiquait plus les jours de communions générales. Ceux qui en avaient fait partie se considéraient
% Source : PDF original, vue 116, côté droit, page manuscrite 233.
toujours comme congréganistes, mais un titre de ce genre qui n'est accompagné d'aucune pratique qui en dérive comme de son principe se réduit à bien peu de choses ; cette dissolution était malheureusement trop réelle et fut bientôt entièrement consommée. Nous verrons plus tard cette Congrégation faire, mais en vain, quelques efforts pour se relever.

La Congrégation des dames subsistait toujours, mais cette société avait cessé de pouvoir être considérée comme une Congrégation. Après le départ du P. Roger elle éprouva sinon une dissolution, du moins une transformation qui lui donna une autre sorte d'existence que celle qu'elle avait eue jusqu'alors. Nous avons dit ailleurs que cette Congrégation se composait comme celle des jeunes gens de trois sections qui portaient les mêmes dénominations et avaient les mêmes attributions. La section des prisons fut obligée de suspendre l'œuvre qu'elle y exerçait ; elle ne voulut pas entreprendre d'aller visiter les malades de l'Hôtel-Dieu, où la Congrégation des demoiselles était déjà en possession de visiter les personnes de leur sexe ; elle ne sut imaginer aucun moyen pour suppléer à celle qui lui échappait, les dames de cette section se trouvèrent donc sans emploi. Depuis la fondation des Congrégations on avait formé dans presque toutes les paroisses de Lyon, des réunions de miséricorde. Les dames qui composaient la section des aumônes faisaient pour la plupart partie de ces sociétés de charité, confondant entre elles l'une et l'autre œuvre. Elles trouvèrent que leur section n'était qu'un double emploi qui devenait superflu ; elles cessèrent donc ce service particulier qu'elles regardèrent comme inutile. Le mal vint essentiellement de ce qu'on ne considéra
% Source : PDF original, vue 117, côté gauche, page manuscrite 234.
pas assez la Congrégation dans ses rapports généraux avec la gloire de Dieu et la sanctification de ses membres et qu'on ne voulut y apercevoir que l'intérêt de ses œuvres, sous ce rapport il est certain que pendant que les deux autres sections étaient réduites à la plus complète inaction, la troisième celle de l'instruction réclamait toutes les ressources disponibles et la coopération de toutes les dames qui voudraient encore aider cette œuvre, non seulement avec leur bourse, mais par un concours actif. Cette section avait pris un très grand développement, elle avait formé pour l'éducation des jeunes filles un grand établissement qui exigeait pour le soutenir les plus grands sacrifices. Dans la disposition où se trouvaient les esprits, cette nouvelle fondation eut bientôt tout envahi, il ne resta alors pour souvenir de la société fondée par le P. Roger, que des assemblées où on se réunit encore tous les mois pour entendre un discours pieux et quelques communions générales. Le nom même de Congrégation a disparu pour faire place à celui de l'œuvre fondée par une de ses sections, et depuis lors elle ne fut plus connue que par le nom qu'elle porte encore aujourd'hui de Providence de Bellecour.

\section*{Transformation de la Congrégation des dames}
% Mention marginale de la page manuscrite 234 ; placée avant la comparaison finale.
La Congrégation des demoiselles avait parfaitement conservé son esprit primitif, elle faisait beaucoup de bonnes œuvres, mais elle avait toujours eu le soin d'opérer dans le secret. Elle rendit encore plus épais le voile dont elle était environnée, et persévéra dans l'union et l'esprit de foi. Aussi continua-t-elle de marcher d'un pas ferme dans sa carrière, et quand il plut à la Providence de nous rendre un peu plus de liberté elle s'était maintenue et se trouva debout.

% Source : PDF original, vue 117, côté droit, page manuscrite 235.
\section*{Marche de la Congrégation}
% Mention marginale des pages manuscrites 235–236.
Revenons maintenant à la Congrégation des jeunes gens. La voilà donc censée dissoute, dès lors, il ne fut plus question ni d'assemblées générales, ni d'élections, ni de réceptions ; mais ses membres étaient plus que jamais unis entre eux et soumis au gouvernement de la Congrégation. Ils apportaient d'autant plus d'exactitude aux assemblées de section que c'était le seul moyen qui leur avait été laissé de se retrouver encore avec plusieurs de leurs frères ; on y mettait toutefois beaucoup de prudence ; toutes les œuvres qui exigeaient un peu de publicité furent abandonnées. Nous avons déjà dit qu'on avait cessé d'aller aux prisons ; on continua quelques visites à l'hôpital mais en se bornant à voir un petit nombre de malades. La section de l'instruction suspendit les catéchismes de persévérance et se contenta de donner ses soins à quelques instructions particulières. Celle des aumônes se partagea la liste de ses pauvres ; chaque membre fut chargé de voir seul ceux qui lui furent assignés.

La Congrégation ne conserva d'autre lien habituel pour unir ensemble les membres des trois sections, que les visites fréquentes qu'y faisaient le préfet et les assistants. On trouvait bien encore à chaque fête le moyen de se réunir auprès des autels, de prier ensemble et de renouveller par la participation à la Ste Communion l'union de tous les cœurs. Rien n'est plus propre à réchauffer le zèle et à augmenter la ferveur qu'un peu de persécution.

Souvent on allait à Fourvières supplier notre auguste patronne de jeter un regard de compassion sur l'Église, sur la France, sur notre Congrégation et d'obtenir enfin pour nous
% Source : PDF original, vue 118, côté gauche, page manuscrite 236, avant le chapitre dixième.
de son divin fils grâce, pardon et miséricorde.

Ces prières étaient entendues, un jour elles seront exaucées, mais auparavant de nouvelles épreuves nous étaient réservées ; elles devaient répandre sur la Congrégation un nouvel éclat par le courage héroïque de ceux de ses membres que le Seigneur avait choisis dans sa miséricorde pour marcher à la suite du chef de l'Église dans les voies de la persécution.

C'est ainsi que Dieu dans sa bonté daignait associer notre Congrégation à ces tribulations révérées par les âmes justes comme des signes de prédestination. L'Esprit-Saint semblait nous faire entendre alors ces paroles adressées autrefois dans l'Apocalypse à l'ange de l'église de Smyrne : Congréganistes de Lyon, ne craignez rien de ce que vous allez avoir à souffrir, car voilà que le démon va jeter quelques-uns d'entre vous dans les prisons, afin que vous soyez mis à l'épreuve ; votre tribulation ne sera que pour un temps ; soyez fidèles et je vous donnerai la couronne de vie.
